% Templat TA TIF UIN Suska Riau -- (c) 2026 Benny Sukma Negara
\chapter{Pendahuluan}
\label{bab:pendahuluan}

%==========================================================
\section{Latar Belakang}
\label{sec:latar}

Latar belakang merupakan subbab yang menjelaskan alasan dan permasalahan yang dihadapi sehingga menghasilkan motivasi untuk melakukan penelitian. \isi{Hapus paragraf petunjuk ini setelah menulis.}

Adapun isi dari subbab Latar Belakang sebaiknya menjelaskan komponen-komponen berikut ini \parencite{lecun2015deep}:
\begin{enumerate}
  \item Permasalahan umum yang dihadapi.
  \item Permasalahan khusus yang dihadapi.
  \item Fenomena-fenomena yang terjadi sebagai akibat dari permasalahan yang terjadi. Sebaiknya fenomena ini diambil dari fakta-fakta riset atau data-data statistik.
  \item Penelitian-penelitian yang sudah dilakukan berkaitan dengan penelitian yang sedang dilakukan peneliti saat ini.
  \item Kelemahan-kelemahan dari penelitian-penelitian sebelumnya, atau aspek penelitian yang belum dilakukan dan yang akan menjadi aspek penelitian oleh peneliti.
  \item Solusi yang ditawarkan, atau penelitian yang akan dilakukan oleh peneliti dan berbagai manfaat dari penelitian tersebut.
\end{enumerate}

Contoh sitasi naratif: menurut \textcite{goodfellow2016deep}, pembelajaran mendalam memungkinkan komputer mempelajari konsep rumit dari konsep yang lebih sederhana.

%==========================================================
\section{Rumusan Masalah}
\label{sec:rumusan}

Rumusan masalah merupakan kalimat dalam bentuk pertanyaan penelitian (\textit{research question}), tidak harus satu kalimat, namun dapat berupa dua atau tiga rumusan masalah. Rumusan masalah harus dibuat sesingkat mungkin.

\isi{Contoh: Bagaimana merancang/membangun ... dan bagaimana kinerjanya diukur menggunakan ...?}

%==========================================================
\section{Batasan Masalah}
\label{sec:batasan}

Batasan masalah digunakan untuk membatasi pembahasan yang akan disampaikan pada laporan penelitian.
\begin{enumerate}
  \item \isi{Data yang digunakan adalah ...}
  \item \isi{Metode/algoritma yang digunakan adalah ...}
\end{enumerate}

%==========================================================
\section{Tujuan Penelitian}
\label{sec:tujuan}

Tujuan penelitian merupakan target yang ingin dicapai dari penelitian ini apabila penelitian ini berjalan dengan baik. \isi{Tujuan sejajar satu-satu dengan rumusan masalah.}

%==========================================================
\section{Manfaat Penelitian}
\label{sec:manfaat}

Manfaat penelitian merupakan penjelasan nilai dan kontribusi yang dihasilkan dari penelitian.
